\documentclass[12pt, a4paper]{article}

% --- Pacchetti Scientifici e Formattazione ---
\usepackage{amsmath, amssymb, amsfonts, amsthm}
\usepackage{booktabs, graphicx}
\usepackage[document]{ragged2e}
\usepackage{xcolor}

% --- Codifica e Font ---
\usepackage[utf8]{inputenc}
\usepackage[T1]{fontenc}
\usepackage[italian]{babel}
\usepackage{mathptmx}

% --- Gestione Margini e Layout ---
\usepackage[includeheadfoot]{geometry}
\geometry{
    a4paper,
    top=3cm,
    bottom=3cm,
    left=3.5cm,
    right=3.5cm,
}

% --- Gestione Interlinea e Figure ---
\usepackage{setspace}
\setstretch{1.3}
\usepackage[figuresright]{rotating}
\usepackage{float}
\usepackage[pdftex, hidelinks]{hyperref}
% Document metadata

\title{Progettazione database}
\author{Mattia Cavina}
\date{03 febbraio 2026}

\begin{document}

\subsection{Importazione del database}

Il database Access attuale è stato convertito e importato direttamente in SQL Server tramite SQL Server Migration Assistant (SSMA)~\cite{SQLServerMigrationAssistant}.
Le tabelle importate sono state successivamente elaborate tramite SQL Server Integration Services (SSIS)~\cite{SQLServerIntegration} per svuotarle e ricrearne la struttura, ottenendo una base di dati pulita, pronta per il nuovo sistema.

\subsection{Normalizzazione e modifica}

Dopo l'importazione sono state effettuate le modifiche e la normalizzazione del database per adattarlo alle esigenze del nuovo configuratore commerciale.

\subsubsection{Convenzioni di denominazione}

Per garantire coerenza e facilità di comprensione, sono state adottate le seguenti convenzioni:

\begin{itemize}
    \item Tabelle e campi in minuscolo; nomi composti separati con underscore (es.\ nome\_tabella).
    \item Nomi dei campi preceduti dal prefisso corrispondente alle prime tre lettere della tabella (es.\ cli\_nome per la tabella clienti).
    \item Per tabelle con nomi composti da due elementi: la prima lettera della prima parola + le prime due lettere della seconda (es.\ cit\_nome per la tabella clienti\_italia).
    \item Per tabelle con nomi composti da tre o più elementi: la prima lettera delle prime tre parole (es.\ cee\_nome per clienti\_extra\_europa).
    \item Chiavi primarie composte da prefisso + id (es.\ cli\_id per la tabella clienti).
    \item Campi booleani con suffisso \_flag dopo il prefisso (es.\ cli\_flag\_sconto).
\end{itemize}

\subsubsection{Campi comuni}
Per facilitare la lettura dei dati e permettere la gestione delle modifiche dei record senza eliminarli ma revisionandoli, si è deciso di aggiungere a tutte le tabelle del database
i seguenti campi:
\begin{itemize}
    \item \textbf{creation\_date} di tipo datetime con default NULL:\ deve essere valorizzata automaticamente con data e ora di creazione del record.
    \item \textbf{last\_edit\_date} di tipo datetime con default NULL:\ deve essere valorizzata automaticamente con data e ora dell'ultima modifica del record.
    \item \textbf{creation\_user} di tipo int con default NULL:\ deve essere valorizzato automaticamente con l'ID dell'utente che ha creato il record.
    \item \textbf{last\_edit\_user} di tipo int con default NULL:\ deve essere valorizzato automaticamente con l'ID dell'utente che ha effettuato l'ultima modifica al record.
    \item \textbf{deleted} di tipo bit con default 0: indica se un record è stato cancellato (cancellazione logica); questa scelta permette di mantenere lo storico dei dati.
\end{itemize}

\subsubsection{Normalizzazione}
Per garantire una maggiore efficienza e coerenza dei dati, si è deciso di normalizzare il database almeno fino alla seconda forma normale (2NF).

\subsection{Non proliferazione tabelle}
\subsubsection{Scopo}
Data l'elevato numero di casistiche rilevate nelle interviste e considerando sviluppi futuri non prevedibili, si è deciso di adottare una strategia di non proliferazione delle tabelle.
Si evita quindi di creare tabelle specifiche per ogni casistica rilevata, preferendo tabelle generiche in grado di contenere più casistiche, in modo da garantire maggiore flessibilità e scalabilità del database.
\subsubsection{Realizzazione}
Per realizzare ciò si è predisposto, in determinate tabelle (es.\ dbo.offerte), un campo ``tipo'' e un campo ``padre'': il primo permette di distinguere la casistica del record (es.\ offerta, ordine, distinta tecnica) e il secondo di ricostruire la gerarchia tra i record (es.\ a quale offerta appartiene un ordine).
Il campo `tipo` è di tipo `int` e, tramite una seconda tabella di lookup (es.\ dbo.tipo\_offerta), è possibile associare a ogni valore del campo `tipo` una descrizione (es. 1 = offerta, 2 = ordine, 3 = distinta tecnica).
Questa soluzione per quanto possa sembrare meno efficiente in termini di prestazioni, garantisce una maggiore flessibilità e scalabilità del database,
permettendo di gestire facilmente nuove casistiche senza dover creare nuove tabelle.
La scelta di adottare questa strategia è stata fatta in particolare per la tabella dbo.offerte dove tramite dei campi attributo possiamo determinare e ricostruire le diverse viste dalla stessa tabella (es: offerta, ordine, distinta tecnica),
evitando di creare tabelle specifiche per ogni casistica rilevata.
Questa scelta è poi applicabile anche all'anagrafica clienti, dove, tramite appositi campi attributo, possiamo determinare e ricostruire se un record è contemporaneamente un cliente, un fornitore o una destinazione.
\subsubsection{Esempio}
\begin{table}[h!]
    \centering
    \begin{tabular}{|c|l|c|l|c|}
        \hline
        \textbf{Id} & \textbf{Codice\_ordine} & \textbf{tipo} & \textbf{codice} & \textbf{ID\_padre} \\ \hline
        1           & OFF14500                & 1             & PRO00001        & null               \\ \hline
        2           & OFF14500                & 1             & PRO00002        & null               \\ \hline
        3           & ORD4850                 & 2             & PRO00001        & 1                  \\ \hline
        4           & ORD4850                 & 2             & PRO00002        & 2                  \\ \hline
        5           & TEC4850                 & 3             & PRO00001        & 3                  \\ \hline
        6           & TEC4850                 & 3             & PRO00002        & 4                  \\ \hline
    \end{tabular}
    \vspace{0.5cm}
    \caption{Tabella Ordini}
\end{table}

\vspace{1cm} % Spazio tra le due tabelle

\begin{table}[h!]
    \centering

    \begin{tabular}{|c|l|}
        \hline
        \textbf{Id} & \textbf{tipo} \\ \hline
        1           & OFFERTA       \\ \hline
        2           & ORDINE        \\ \hline
        3           & TECNICA       \\ \hline
    \end{tabular}
    \vspace{0.5cm}
    \caption{Tabella Decodifica}
\end{table}
\subsection{Gestione storico dei dati}
Per la gestione dello storico dei dati si è scelto di predisporre tabelle di backup che contengano i dati storici, in modo da non appesantire le tabelle principali.
Questa soluzione è supportata da trigger che spostano automaticamente i dati storici nelle tabelle di backup, garantendo una gestione efficiente dello storico.
\subsection{Pulizia database importato}
A seguito dell'importazione del database, è stato necessario effettuare una pulizia delle tabelle, eliminando tabelle e campi non più necessari.
L'analisi è stata effettuata mediante interviste con gli utenti e tramite l'analisi dei dati, valutando se fossero ancora utilizzati.
Di seguito sono riportati i passi seguiti per la pulizia del database importato.


\subsubsection{Accorpamento tabelle}

\paragraph{Agenti, autori\_offerta}
Le tabelle degli agenti e degli autori\_offerta potrebbero essere accorpate in un'unica tabella, differenziando i ruoli tramite un attributo.
Le due tabelle sono molto simili e condividono molti campi, quindi potrebbe essere più efficiente gestirle in un'unica struttura.

\paragraph{lingue\_accessori, lingue\_documenti}
Sono entrambe tabelle che contengono le lingue disponibili per gli accessori e per i documenti; potrebbero essere accorpate in un'unica tabella per la scelta della lingua, gestendo le traduzioni esternamente al database.

\paragraph{offerte\_dettagli, offerte\_dettagli\_optional, offerte\_ra\_dettagli}
Queste tabelle sono molto simili e hanno scopi e strutture affini; potrebbero essere accorpate in un'unica tabella con un attributo che indichi se si tratta di un prodotto opzionale.

\paragraph{offerte\_ra, offerte}
Queste tabelle sono molto simili e condividono molti campi; potrebbero essere accorpate in un'unica tabella, differenziando il tipo di offerta tramite un attributo.

\paragraph{clienti, utilizzatori}
Queste tabelle condividono dei campi e talvolta il cliente è anche l'utilizzatorefinale.
Potrebbero essere accorpate in un'unica tabella con un attributo che identifica se è un cliente o cliente finale.
Importante creare anche un campo padre in cui il cliente finale può essere collegato al cliente, in modo da poter
gestire le relazioni tra i clienti ed utilizzatori.
\paragraph{tipi\_causali, tipi\_assistenza}
Le tabelle hanno il medesimo scopo di gestire le causali di vendita;
potrebbero essere accorpate in un'unica tabella con un attributo che identifica il tipo di causale.


\subsubsection{Eliminazione tabelle}
Di seguito sono elencate le tabelle che potrebbero essere eliminate in quanto non più necessarie o ridondanti:
\begin{itemize}
    \item agenti
    \item associazioni\_ted
    \item lingue\_accessori
    \item lingue\_documenti
    \item listini
    \item offerte\_dettagli\_optional
    \item offerte\_ra
    \item offerte\_ra\_dettagli
    \item ordini (Con dovute accortezze.)
    \item prodotti\_storico\_prezzi
    \item riferimenti\_esterni
    \item riferimenti\_interni
    \item tipi\_cliente
    \item utilizzatori
\end{itemize}

\subsection{Tabelle da implementare}
In base alle interviste, è necessario implementare ulteriori tabelle, tra cui:
\begin{itemize}
    \item \textbf{Tabelle dati storici}: tabelle per la gestione dello storico e delle revisioni dei prodotti, in modo da conservare lo storico dei prezzi e delle caratteristiche senza eliminare i prodotti obsoleti.
          Questa tabella sarà gestita tramite trigger che sposteranno automaticamente i dati obsoleti nella tabella di backup; la ricostruzione avverrà poi tramite query basate sulla data di validità.
    \item \textbf{Tabelle macro prodotti}: tabella centrale per migliorare la qualità delle offerte.
          Permette di creare macro-categorie per la compatibilità dei prodotti; scegliendo la macro-categoria verranno suggeriti solo i prodotti compatibili, migliorando la qualità delle offerte e riducendo gli errori.
          In questa tabella dovranno essere associate anche le proprietà, analogamente ai prodotti, per gestire le proprietà anche a livello di macro-categoria.
    \item \textbf{Tabella notifiche}: per la gestione delle notifiche, con campi per utente e data, in modo da filtrare le notifiche non lette.
    \item \textbf{Tabella feedback}: per la raccolta e la gestione dei feedback, con vista amministratore dedicata per l'analisi e il miglioramento del sistema.
    \item \textbf{Tabella listino merceologico}: permette di applicare sconti o maggiorazioni in base al gruppo merceologico e alla tipologia di prodotto.
    \item \textbf{Tabella proprietà prodotti}: per gestire le proprietà dei prodotti in modo dinamico e strutturato; sono necessarie tre tabelle in relazione molti a molti.
          La prima (`tipi\_proprieta`) conterrà le tipologie di proprietà (es. meccaniche, elettriche, pneumatiche, software, ecc.).
          \begin{table}[h!]
              \centering
              \begin{tabular}{|c|l|}
                  \hline
                  \textbf{Id} & \textbf{Tipo\_proprieta} \\ \hline
                  1           & Meccaniche               \\ \hline
                  2           & Elettriche               \\ \hline
                  3           & Pneumatiche              \\ \hline
                  4           & Software                 \\ \hline
              \end{tabular}
              \vspace{0.5cm}
              \caption{Tabella Tipi Proprietà}
          \end{table}
          La seconda (`proprieta`) e la relativa tipologia (es. peso, potenza, tensione, ecc.) con una colonna `tipo\_proprieta` per associarla alla tipologia.
          \begin{table}[h!]
              \centering
              \begin{tabular}{|c|l|c|}
                  \hline
                  \textbf{Id} & \textbf{Proprieta} & \textbf{Tipo\_proprieta} \\ \hline
                  1           & Peso               & 1                        \\ \hline
                  2           & Potenza            & 1                        \\ \hline
                  3           & Tensione           & 2                        \\ \hline
                  4           & Portata            & 3                        \\ \hline
              \end{tabular}
              \vspace{0.5cm}
              \caption{Tabella Proprietà}
          \end{table}
          La terza (`proprieta\_prodotto`) sarà una tabella di associazione tra prodotti e proprietà in relazione 1 a molti, questa tabella conterrà i campi `prodotto\_id`,
          `nome\_proprieta\_id` e `valore\_proprieta` per associare ogni prodotto alle sue proprietà specifiche.
          \begin{table}[H]
              \centering
              \begin{tabular}{|c|c|c|}
                  \hline
                  \textbf{Prodotto\_id} & \textbf{Nome\_proprieta\_id} & \textbf{Valore\_proprieta} \\ \hline
                  1                     & 1                            & 100 kg                     \\ \hline
                  1                     & 2                            & 2000 W                     \\ \hline
                  2                     & 3                            & 220 V                      \\ \hline
                  3                     & 4                            & 500 l/min                  \\ \hline
              \end{tabular}
              \vspace{0.5cm}
              \caption{Tabella Proprietà Prodotto}
              Questa struttura ci permette di aggiungere nuove proprietà dal software senza dover modificare la struttura del database.
          \end{table}
    \item \textbf{Tabella capitoli offerte}: tabella per gestire delle sottosezioni all'interno delle offerte, in modo da organizzare meglio le offerte e permettere di gestire meglio le offerte complesse.
          Questa tabella sarà collegata alla tabella `offerte` tramite una chiave esterna, permettendo di associare più capitoli a un'offerta.
          L'uso di una tabella ci permette anche di riordinare i capitoli in modo dinamico, senza dover modificare la struttura dell'offerta.
    \item \textbf{Tabella destinazioni impianto}: tabella dedita a raccogliere tutte le destinazioni di impianto, così da non doverle riscrivere per cliente, utilizzatore o destinazione. Viene così eliminata la destinazione dalla gestione inizialmente pensata nella tabella clienti.
          La tabella sarà collegata alla tabella `clienti` tramite una chiave esterna, permettendo di associare più destinazioni a un cliente.
    \item \textbf{Tabella stato offerta}: tabella mancante di lookup per la gestione dello stato di offerta e ordine (es.\ preventivo, in approvazione, approvato, rifiutato; in lavorazione, spedito, consegnato).
          Questo valore andrà inserito come chiave esterna nella tabella `offerte`, per permettere la gestione dello stato.
    \item \textbf{Tabelle di lookup}: per la gestione centralizzata delle decodifiche delle foreign key dei campi dediti alla definizione di attributi di una tabella (es. tipologia di offerta, tipologia di cliente, stato dell'offerta, ecc.).
          L'approccio ci permette di getsire come chiave esterna un intero numero, evitando di dover gestire stringhe e garantendo una maggiore flessibilità e scalabilità del database.
          Queste tabelle si distinguono per il prefisso ``tipi\_''.
          Le tabelle necesarie sono:
          \begin{itemize}
              \item \textbf{Tipi causali}: per la gestione delle causali di vendita (es. VSS \- vendita, RIT \- reso, ecc.).
              \item \textbf{Tipi cliente}: per la gestione delle tipologie di cliente(es.\ cliente finale, rivenditore, ecc.).
              \item \textbf{Tipi documento}: per la gestione delle tipologie di documento (es. offerta, ordine, distinta tecnica).
              \item \textbf{Tipi imballo}: per la gestione delle tipologie di imballo per le spedizioni (es.\ su gomma, via nave, ecc.).
              \item \textbf{Tipi impianto}: per la gestione delle tipologie di impianto (es. Bakery, Biscuit, Chocolate,ecc.).
              \item \textbf{Tipi pagamento}: per la gestione delle tipologie di condizioni di pagamento (es. 30 giorni, 60 giorni, 90 giorni, ecc.).
              \item \textbf{Tipi percetuale agente}: per la gestione delle tipologie di percentuale agente (es. 1\%, 2\%, 3\%, ecc.).
              \item \textbf{Tipi proprietà}: per la gestione delle tipologie di proprietà dei prodotti (es. meccaniche, elettriche, pneumatiche, software, ecc.).
              \item \textbf{Tipi resa}: per la gestione delle tipologie di resa ovvero gli INCOTERMS (es.\ franco fabbrica, franco destino, ecc.).
              \item \textbf{Tipi stato offerta}: per la gestione delle tipologie di stato dell'offerta e dell'ordine (es.\ preventivo, in approvazione, approvato, rifiutato; in lavorazione, spedito, consegnato).
              \item \textbf{Tipi vendita}: per la gestione delle tipologie di vendita (es.\ assistenza, ricambi, vendita).
              \item \textbf{Tipi aree contropartita}: per la gestione delle tipologie di aree contropartita (es. Italia, Europa, Extra Europa).
          \end{itemize}
    \item \textbf{Tabella gestione permessi}: tabella contenente l'elenco delle tabelle del database e i permessi per ciascuna di lettura,scrittura e cancellazione, in modo da gestire centralmente i permessi di accesso alle tabelle.
          Questa tabella sarà collegata alla tabella `utenti` tramite una chiave esterna, permettendo di associare i permessi in base alla tipologia di utente.
\end{itemize}

\subsection{Seconda Intervista}
Una volta definita la struttura del database, è stata effettuata una seconda intervista con gli utenti per validare
la struttura e raccogliere eventuali feedback o modifiche da apportare.
La prima intervista si è svolta con la CEO dell'azienda che ci ha chiarito i dubbi riguardanti le tabelle non più necessarie e
ci ha confermato la scelta di adottare una strategia di non proliferazione delle tabelle.
Una aspettoa ccolto positivamente è la scelta di gestire più tipologie di documenti per fare le valutazioni tecniche post vendita,
in modo da minimizzare gli errori e migliorare la qualità delle offerte.
Durante la stessa riunione è definita la roadmap per la realizzazione del software.

\subsubsection{Allineamneto con i commerciali}
Corretto il database con le modifiche raccolte dalla intervista si è fatto un incontro con il resposnsabile commerciale che
ci ha appoggiato nelle scelte, rimarcando l'importanza di migliorare la parte di offerte di ricambistica.

L'ultimo incontro si è svolto con due commerciali di riferimento che hanno fatto emergere alcuni dubbi e richieste di modifica come:
\begin{itemize}
    \item Possibiltà di dare un titolo all'offerta.
    \item Gestione degli sconti e rincari non era ottimale nella prima versione , si sono quindi aggiunti i campi necessari.
    \item Hanno posto l'interrogativo sulla gestione del "Caso speciale", ovvero quando i clienti richiedono prodotti non presenti a catalogo.
    \item Dall'intrevista è emerso il problema della ricerca dei macrogruppi di riferimento dei prodotti, si è quindi deciso di aggiungere una ulteriore tabella proprieta\_macro\_gruppi per catalogarli meglio.
    \item Da parte del commerciale addetto ai ricambi è stata avanzata la richeista di un campo per salavre il numero di tiket di assistenza, in modo da poter collegare l'offerta di ricambio al ticket di assistenza e migliorare la tracciabilità delle offerte di ricambio.
\end{itemize}

\subsection{Definizione tabelle backup}
La gestione dello storico dei dati verrà fatto mediante tabelle dedicate di backup, ma non sarà necessario tenere traccia di tutte le tabelle,
ma solo di quelle più importanti e critiche per il funzionamento del sistema.
L'automatismo di spostamento dei dati obsoleti nelle tabelle di backup sarà gestito tramite trigger,
in modo da garantire una gestione efficiente dello storico senza appesantire le tabelle principali.
La scelta del trigger ha un doppio aspetto da un lato permette di gestire in modo automatico lo spostamento dei dati obsoleti,
dall'altro lato potrebbe sia appesantire le prestazioni che complicare la sincronizazione dei dati che verranno importati da database locali.
Si è scelto di addottare comunque questa soluzione in quanto semplifica la gestione lato sofware.
Le tabelle soggette a questa gestione sono:
\begin{itemize}
    \item \textbf{Clienti} Tabella dei clienti ed utilizzatori.
    \item \textbf{Clienti\_destinazioni} Tabella delle destinazioni dei clienti.
    \item \textbf{Documenti} Tabella dei documenti di offerta,ordini,tecniche.
    \item \textbf{Documenti\_righe} Tabella delle righe dei documenti.
    \item \textbf{Montaggi} Tabella dela costificazione dei montaggi.
    \item \textbf{Prodotti} Tabella dei prodotti.
    \item \textbf{Prodotti\_composizioni} Tabella delle composizioni dei prodotti, ovvero la loro distinta.
    \item \textbf{Proprietà\_valori} Tabella dei valori delle proprietà del prodotto.
\end{itemize}

\subsection{Popolamento delle tabelle}
Raggiunto un buon grado di completezza, si è proceduto al popolamento delle tabelle con i dati necessari per il funzionamento del sistema.

Come per la costruzione il popolamento è avvenuto tramite SSIS dove si sono creati diversi container per suddividere il lavoro in più fasi,
partendo dalle tabelle di lookup, passando poi alle tabelle dove i dati sono presi dal vecchio sistema ed infine un container per le tabelle dove presenti i dati fittizzi.

La divisione del lavoro in più container ci ha permesso di gestire meglio il processo di popolamento, oltre che in una fase futura ci permetterà di 
disabilitare il container di popolamento dei dati fittizzi, in modo da non sovrascrivere i dati reali con quelli fittizzi.
Per quanto riguarda invece le tabeelle importate verranno controllate in un secondo momento, in modo da eliminare eventuali dati non più necessari o obsoleti,
e per verificare che i dati siano stati importati correttamente.

\section{Popolamento database con dati reali}
Raggiunto un buon grado di completezza, si è proceduto al popolamento delle tabelle con i dati reali 
necessari per il funzionamento del sistema.
Il popolamento come scelto è avvenuto prendeno i dati dai programmi adibiti a Master Data come ad esempio:
\begin{itemize}
    \item \textbf{AD-HOC}: per i dati anagrafici dei clienti, utilizzatori e prodotti.
    \item \textbf{Fusion}: per i dati tecnici dei prodotti, come le caratteristiche e le proprietà.
\end{itemize}
\subsection{Struttura dati aziendale}
Tutti i dati che necessitiamo risiedono in un unico server aziendale, ma sono gestiti da diversi programmi,
 ognuno con il proprio database e la propria struttura dati, ma tutti con dtaabse SQL Server.
  Per poterli importare nel nostro sistema è necessario fare delle query di estrazione dei dati, esse 
  sono state create facendo reverser engineering dei singoli databse.
  Una volta capita la struttura dei dati, si è proceduto a creare delle query di estrazione dei dati,
   in modo da poterli importare nel nostro sistema.

\subsection{Approcio scelto}
Si è scelto di creare delle stored procedure che ricercano ed eventualmente ricostruiscono i dati in base alle politiche del
nostro sistema, in modo da avere un controllo maggiore sui dati importati e garantire la coerenza dei dati.
Le stored procedure sono state scelte in quanto le query di importazioni sono da eseguire in modo automatico e periodico,
quindi è necessario avere un meccanismo di controllo facilmente gestibile.
Come per le tabelle anchesse sono integrate nella creazione del database in SSIS.

\subsection{Ordine di creazione}
Le stored procedure create hanno un ordine di esecuzione ben definito, in quanto alcune dipendono da altre per la creazione dei dati.
L'ordine di esecuzione è il seguente:
\begin{enumerate}
    \item \textbf{sp\_import\_prodotti}: importa i prodotti e le loro caratteristiche.
    \item \textbf{sp\_import\_distinte\_prodotti}: importa le distinte dei prodotti e le loro caratteristiche.
    \item \textbf{sp\_import\_proprieta}: importa le proprietà dei prodotti.
    \item \textbf{sp\_import\_proprieta\_valori}: importa i valori delle proprietà del prodotto.
\end{enumerate}

Separatamnete:
\begin{enumerate}
    \item \textbf{sp\_import\_clienti}: importa i clienti e le loro destinazioni.
    \item \textbf{sp\_import\_utenti}: importa gli utenti del sistema.
\end{enumerate}
\end{document}

